% Source: https://artofproblemsolving.com/community/c3h1716751_algebra_problems_with_n0
% Stable identity: aops-topic-1716751-q01
Doing the same thing to both sides of a true equation produces a true equation: if $5\cdot6=n$, then certainly $(5\cdot6)^0=n^0$. This step is an implication in one direction only,
\[
5\cdot6=n\ \Longrightarrow\ 1=1 ,
\]
and it can be reversed only when the operation applied is one-to-one, so that equal outputs come from equal inputs (for example adding $3$, multiplying by a nonzero number, or taking a cube).

The map $t\mapsto t^0$ sends every nonzero number to $1$, so it is not one-to-one; it throws away all information about $t$. From $n^0=1$ we learn only that $n\ne0$ (as $0^0$ is usually left undefined), not that $n=30$. The final statement ``$1=1$ for every $n$'' says nothing about which $n$ satisfy the \emph{original} equation; the set of solutions of the new equation is larger than that of the old one.

Concretely, $n=7$ satisfies $1=1$ but not $5\cdot6=n$. The only solution of the original equation is $n=\boxed{30}$. (The same mistake appears when both sides are squared: $x=3$ leads to $x^2=9$, which also has the extra solution $x=-3$.)
